\documentclass[10pt,a4paper]{amsart}
\usepackage[margin=19mm]{geometry}
\usepackage{amsmath,amssymb,mathtools}
\usepackage[scr=rsfs,cal=euler]{mathalfa}
\usepackage{xfrac}
\usepackage[unicode,hidelinks,bookmarks=false]{hyperref}
\newcommand{\ie}{i.e.\ }
\newcommand{\cf}{cf.\ }
\newcommand{\resp}{respectively\ }
\newcommand{\Subsection}{\subsection}
\providecommand{\nobreakdash}{\nobreak\hbox{-}}
\setlength{\emergencystretch}{2em}
\multlinegap0pt
\usepackage{amssymb}
\usepackage{enumerate}
\usepackage{xcolor}
\newtheorem{lemma}{Lemma}
\title{Laguerre-Sobolev orthogonal Polynomials and Boundary Value Problems on a semi-infinite domain}
\author{Cleonice F. Bracciali, Miguel A. Pi\~nar}
\date{Source: arXiv:2602.06685v1 (CC BY 4.0)}
\begin{document}
\maketitle

\noindent\emph{Source and licence.} Laguerre-Sobolev orthogonal Polynomials and Boundary Value Problems on a semi-infinite domain. Version arXiv:2602.06685v1 (CC BY 4.0), distributed by arXiv under the Creative Commons Attribution 4.0 International licence, http://creativecommons.org/licenses/by/4.0/, as recorded in the arXiv licence metadata for this exact version and shown on the abstract page https://arxiv.org/abs/2602.06685v1. Full licence notice: \texttt{licenses/arxiv-2602.06685-CC-BY-4.0.txt}.\par\smallskip

\noindent\textit{Editorial scope.} Counted unit: the article's lemma rat\_func\_an (source0.tex lines 343--351). The article's complete original proof of it is delivered with it (source0.tex lines 352--399, one \texttt{proof} environment of 48 lines). No mathematics is written, repaired or re-derived: the statement and the proof are the article's own lines, in the article's own notation, and no retained line is substituted or re-flowed.  Retained uncounted as the source-local context the proof needs: lines 108--113; lines 153--156; lines 244--246; lines 289--291.  Nothing was omitted from any retained span, so the package carries no omission record.  No retained line contains a citation, so the wrapper carries no bibliography and no bibliography entry.  No retained line contains a figure environment, an image-inclusion command or drawing code, so no image is delivered and the package declares no asset dependency.  Editorial formatting only, in addition: an AMS \texttt{amsart} preamble that restates the article's own declarations and macros used by the retained lines, namely the environments \texttt{lemma} on separate counters, together with the standard packages those lines need.  The retained environments print in this document as Lemma 1 in order of appearance, whereas the article prints the same items with the numbering of its own full document.  The complete native source of the article as distributed in the arXiv e-print for this exact version is bundled unchanged as \texttt{sources/194116/source0.tex}.
\begin{equation}\label{ttrr}
\begin{aligned}
& (n+1)\,L_{n+1}^{(\alpha)}(x) = (2n+1 + \alpha -x)L_{n}^{(\alpha)}(x) - (n+\alpha) \,L_{n-1}^{(\alpha)}(x), \quad n\geqslant 0,\\
& L_{-1}^{(\alpha)}(x) = 0, \quad L_0^{(\alpha)}(x) = 1.
\end{aligned}
\end{equation}

\begin{equation} \label{asymp_expan}
\frac{L_{n+j}^{(\alpha)}(z)}{L_{n}^{(\beta)}(z)}
=\left(-\frac{z}{n}\right)^{\frac{\beta-\alpha}{2}}\sum_{m=0}^{d-1} U_m(\alpha,\beta,j,z)n^{-m/2}+\mathcal{O}(n^{-d/2}),
\end{equation}

\begin{equation}\label{conechat}
L_{n}^{(1)} (x)= S_{n}(x) + {a}_{n-1}S_{n-1}(x), \quad n\geqslant 1,
\end{equation}

\begin{equation} \label{rec_rel_an}
a_n = \frac{n+2}{4 \lambda + 2(n+1) - n a_{n-1}}
\end{equation}

\begin{lemma} The sequence $\{a_n\}_{n\geq 0}$ in \eqref{conechat} satisfies
\begin{equation} \label{rat_func_an}
a_n  = \frac{n+2}{n+1} \frac{L^{(1)}_{n}(-4\lambda)}{ L^{(1)}_{n+1}(-4\lambda)}, \quad
\end{equation}
for $n\geqslant 0$, and therefore
\begin{equation} \label{lim_an}
 a_n =1 - \frac{\sqrt{4\lambda}}{\sqrt{n}} +\mathcal{O}(\frac{1}{n}) .
\end{equation}
\end{lemma}

\begin{proof}
First, we consider the sequence of polynomials in $\lambda$, $\{q_n(\lambda)\}_{n\geqslant 0}$, satisfying the recurrence relation
\begin{equation} \label{ttrr_qn}
q_{n+1}(\lambda) = [ 4\lambda +2(n+1)] q_{n}(\lambda) -n(n+1) q_{n-1}(\lambda), \quad n\geqslant 1,
\end{equation}
where  $q_{0}(\lambda) = 1 $ and $q_{1}(\lambda)= 4\lambda +2$.


We see, for $n=0$, that
$$
a_0 = \frac{2}{4 \lambda + 2} = \frac{2 \,q_{0}(\lambda)}{q_{1}(\lambda)}.
$$
Next, we assume
$$
a_{n-1} = \frac{(n+1) q_{n-1}(\lambda)}{q_{n}(\lambda)},
$$
Hence, from \eqref{rec_rel_an}, we have
\begin{align*}
%a_n = &  \frac{n+2}{4 \lambda + 2(n+1) - n a_{n-1}} \\
a_n = &  \frac{n+2}{4 \lambda + 2(n+1) - n  (n+1)\frac{ q_{n-1}(\lambda)}{q_{n}(\lambda)}
} \\
= &  \frac{(n+2)q_{n}(\lambda)}{[4 \lambda + 2(n+1)]q_{n}(\lambda) - n  (n+1) q_{n-1}(\lambda)}.
\end{align*}
and from \eqref{ttrr_qn} we conclude
$$
a_n  = \frac{(n+2) q_{n}(\lambda)}{q_{n+1}(\lambda)}, \quad n\geqslant 0.
$$

The three-term recurrence relation \eqref{ttrr} with $\alpha =1$, becomes
$$ (n+1)\,L_{n+1}^{(1)}(x) = [-x+ 2(n+1) ]L_{n}^{(1)}(x) - (n+1) \,L_{n-1}^{(1)}(x), \quad n\geqslant 1,
$$
with $ L_{0}^{(1)}(x) =1, $ and $ L_{1}^{(1)}(x) = -x+2.$
By comparing  with the three-term recurrence relation \eqref{ttrr_qn}, with $x=-4\lambda$, we obtain
$$
q_n(\lambda) =  n! \, L_{n}^{(1)}(-4\lambda), \quad
n \geqslant 0.
$$
Therefore, for  $n\geqslant 0$,
$$
a_n  = \frac{(n+2) q_{n}(\lambda)}{q_{n+1}(\lambda)}
=  \frac{n+2}{n+1}
\frac{L^{(1)}_{n}(-4\lambda)}{L^{(1)}_{n+1}(-4\lambda)},
$$
and \eqref{rat_func_an} holds.

Finally, \eqref{lim_an} follows from \eqref{rat_func_an} and \eqref{asymp_expan}
with $\alpha = \beta = 1, j = -1,$ and $d=2$.
\end{proof}

\end{document}
